\documentclass[10pt]{article}

\usepackage[T1]{fontenc}
\usepackage[utf8]{inputenc}
\usepackage[english]{babel}
\usepackage[margin=0.72in]{geometry}
\usepackage{graphicx}
\usepackage{xcolor}
\usepackage{microtype}
\usepackage[default,scale=0.95]{sourcesanspro}
\renewcommand{\familydefault}{\sfdefault}
\usepackage[colorlinks=true,urlcolor=blue!50!black]{hyperref}
\usepackage{tabularx}
\usepackage{booktabs}
\usepackage{array}

\definecolor{deepblue}{HTML}{163A59}
\definecolor{accent}{HTML}{D97706}
\definecolor{muted}{HTML}{4B5563}

\setlength{\parindent}{0pt}
\setlength{\parskip}{0.4em}

\newcolumntype{L}{>{\raggedright\arraybackslash}p{0.14\linewidth}}
\newcolumntype{R}{>{\raggedright\arraybackslash}p{0.82\linewidth}}

\begin{document}
\thispagestyle{empty}

{\Large\bfseries\color{deepblue} Agentic AI for Economic Research}\\[-0.15em]
{\large Practical workshop for first-year PhD students in Economics}\\[0.25em]
{\color{accent}\rule{\linewidth}{1.2pt}}

\vspace{0.35em}

\textbf{Format:} 4 sessions $\times$ 2 hours (8 hours total) \\
\textbf{Instructor:} Prof. Marco Le Moglie -- Catholic University of the Sacred Heart

\vspace{0.35em}

\textbf{\color{deepblue}Why this course}

Economic research still runs on the same steps---literature, data, identification,
estimation, writing, revision---but the tools around those steps are changing quickly.
Agents can now read files, run code, search the web, and carry out multi-step tasks on
your behalf. That can save real time, but only if you know how to set them up, what to
delegate, and what to verify.

This course is not about whether AI will replace economists, nor about choosing the
``best'' model. It is a practical introduction to \emph{agentic workflows}: how to use
AI assistants as controlled collaborators across a research pipeline, without losing
oversight of what they do.

\vspace{0.2em}

\textbf{\color{deepblue}Structure}

\textbf{Before class --- Pre-course readings (required).}
Theoretical foundations are covered in written materials distributed before Session~1:
what agents are, how they differ from chatbots, core vocabulary (context, tokens, rules,
MCP connectors), and the main risks (hallucination, automation bias, over-delegation).
Come to the first session having read these notes.

\vspace{0.15em}

\begin{tabularx}{\linewidth}{@{}L R@{}}
\toprule
\textbf{Session} & \textbf{Topic} \\
\midrule
\textbf{1} &
\textbf{Interfaces and setup.} Desktop apps, IDEs, CLIs, and browser-based tools;
chat vs.\ agent modes; model and context settings; project rules; basic permissions
and control. Goal: choose and configure a working environment. \\
\addlinespace
\textbf{2} &
\textbf{Memory and knowledge.} Persistent memory, RAG, and project knowledge bases;
building and querying a research ``second brain'' with Obsidian. Goal: give agents
stable context across sessions without dumping everything into every prompt. \\
\addlinespace
\textbf{3} &
\textbf{Research workflow --- live examples (I).} Concrete cases along the pipeline:
literature synthesis, data documentation, replication checks, Stata/LaTeX assistance,
referee-response drafting. Emphasis on prompting, verification, and iteration. \\
\addlinespace
\textbf{4} &
\textbf{Research workflow --- live examples (II).} Going further: custom commands,
hooks, skills, automations, agent loops, and lightweight multi-agent review (e.g.\
code audit, robustness checks). Goal: reusable setups you can adapt to your own
projects. \\
\bottomrule
\end{tabularx}

\vspace{0.2em}

\textbf{\color{deepblue}What you will be able to do}

By the end of the course, you should be able to:
\begin{enumerate}\setlength\itemsep{0.15em}
  \item \textbf{Set up} an agentic environment suited to economic research (IDE or CLI,
        rules, connectors).
  \item \textbf{Decide} what to delegate to an agent and what to keep under direct
        control.
  \item \textbf{Build} a small knowledge base and connect it to your workflow so agents
        can retrieve project-specific context.
  \item \textbf{Apply} agents to common research tasks---reading papers, documenting
        data, coding, writing, responding to feedback---with explicit verification
        steps.
  \item \textbf{Design} simple reusable automations (commands, skills, hooks) for tasks
        you repeat often.
  \item \textbf{Evaluate} agent output critically: spot errors, avoid automation bias,
        and maintain an audit trail.
\end{enumerate}

\vspace{0.15em}

\textbf{\color{deepblue}What this course is not}
Not a programming course (basic familiarity with Stata/R/Python and LaTeX is enough);
not a survey of every AI product on the market; not a substitute for econometric
training or research judgment.

\vfill
{\scriptsize\color{muted}
Pre-course readings, slide sources, and session recordings:
\href{https://github.com/MarcoLeMoglie/Agentic_Economic_Research}{github.com/MarcoLeMoglie/Agentic\_Economic\_Research}.
}

\end{document}
